\chapter{Introduction to File Handling in Python (Revisited)}

\section{Introduction: Real-Life Hook}

Think about tasks students often do:

\begin{itemize}
    \item Keeping a journal or diary  
    \item Saving homework assignments or notes  
    \item Storing lists of contacts, scores, or favorite books  
\end{itemize}

Python lets us \textbf{store data in files}, read it later, and update it. This is called \textbf{file handling}.

\section{Concept: What is File Handling?}

\begin{definition}{File Handling}
File handling is the process of \textbf{creating, reading, writing, and updating files} using a programming language.
\end{definition}

- Files store information \textbf{permanently on a disk}  
- Two common types:  
  \begin{itemize}
    \item Text files (\texttt{.txt}) — simple readable content  
    \item Binary files (\texttt{.bin}, \texttt{.jpg}, \texttt{.png}) — non-text files like images, videos, etc.  
  \end{itemize}

\section{Opening and Closing Files}

\begin{verbatim}
# Open a file for reading
file = open("notes.txt", "r")
content = file.read()
print(content)
file.close()

# Open a file for writing
file = open("notes.txt", "w")
file.write("This is my first note.\n")
file.close()
\end{verbatim}

- \texttt{"r"} = read mode  
- \texttt{"w"} = write mode (overwrites existing content)  
- \texttt{"a"} = append mode (adds content at the end)  
- Always \textbf{close files} after use  

\section{Using the with Statement (Recommended)}

\begin{verbatim}
with open("notes.txt", "r") as file:
    content = file.read()
    print(content)

with open("notes.txt", "a") as file:
    file.write("Adding another note.\n")
\end{verbatim}

- \texttt{with} automatically \textbf{closes the file}  

\section{Reading Files Line by Line}

\begin{verbatim}
with open("notes.txt", "r") as file:
    for line in file:
        print(line.strip())
\end{verbatim}

- \texttt{strip()} removes extra spaces and newlines  

\section{Mini-Project: Student Score Logger}

Concept:

\begin{itemize}
    \item Save student names and marks in a file  
    \item Allow adding new entries  
    \item Read all entries and calculate average marks  
\end{itemize}

Example:

\begin{verbatim}
# Writing student scores
with open("scores.txt", "a") as file:
    name = input("Enter student name: ")
    marks = input("Enter marks: ")
    file.write(f"{name},{marks}\n")

# Reading student scores
total = 0
count = 0
with open("scores.txt", "r") as file:
    for line in file:
        name, marks = line.strip().split(",")
        print(f"{name}: {marks}")
        total += int(marks)
        count += 1

average = total / count
print(f"Average Marks: {average}")
\end{verbatim}

- Students learn \textbf{reading, writing, parsing, and calculating} from files  

\section{Step-by-Step Diagram: File Handling Workflow}

\begin{center}
\begin{tikzpicture}[node distance=1.5cm]
\node (start) [draw, ellipse, fill=green!20] {Start Program};
\node (open) [draw, rectangle, fill=yellow!20, below of=start] {Open File (r, w, a)};
\node (process) [draw, rectangle, fill=orange!20, below of=open] {Read or Write Data};
\node (use) [draw, rectangle, fill=yellow!20, below of=process] {Process or Display Data};
\node (close) [draw, ellipse, fill=green!20, below of=use] {Close File};
\draw[->] (start) -- (open);
\draw[->] (open) -- (process);
\draw[->] (process) -- (use);
\draw[->] (use) -- (close);
\end{tikzpicture}
\end{center}

\section{Key Concepts}

\begin{itemize}
    \item Files store \textbf{permanent data}  
    \item Modes: \texttt{"r"} (read), \texttt{"w"} (write), \texttt{"a"} (append)  
    \item Always close files or use \texttt{with} statement  
    \item Can read files line by line or as a whole  
    \item Parsing data allows \textbf{calculations, searches, and reports}  
\end{itemize}

\section{End-of-Chapter Exercises}

\begin{exercise}{Very Short Questions}
\begin{enumerate}
    \item What is file handling?  
    \item Name three file modes.  
    \item What is the advantage of using \texttt{with} statement?  
\end{enumerate}
\end{exercise}

\begin{exercise}{Short Answer Questions}
\begin{enumerate}
    \item Write Python code to read a file line by line and print each line.  
    \item Explain the difference between \texttt{"w"} and \texttt{"a"} modes.  
    \item How do you split a CSV line into name and marks?  
\end{enumerate}
\end{exercise}

\begin{exercise}{Activity-Based Questions}
\begin{enumerate}
    \item Create a program that stores your favorite movies in a file and reads them back.  
    \item Build a mini-address book that allows adding names and phone numbers and displaying all contacts.  
    \item Modify the Student Score Logger to save data in two separate files: one for names and one for marks.  
\end{enumerate}
\end{exercise}

\section{Teacher Notes}

\marginnote{
\textbf{Learning Objectives:}  
\begin{itemize}
    \item Students can create, read, write, and append files  
    \item Students understand file modes and best practices  
    \item Students can process data from files for calculations and reports  
\end{itemize}

\textbf{Teaching Tips:}  
\begin{itemize}
    \item Start with text files before introducing CSV or structured formats  
    \item Demonstrate reading and writing step by step  
    \item Encourage students to experiment with real-life examples like notes, scores, or contacts  
\end{itemize}

\textbf{Assessment:}  
\begin{itemize}
    \item Observe correct file creation and reading/writing  
    \item Check parsing of data and calculations  
\end{itemize}
}